\chapter{Clearing and Settlement}\label{ch:m1:clearing-and-settlement}

At 5:11 on the morning of 28 January 2021 an automated notice arrived at a
retail broker's \omterm{def:m1:clearing-and-settlement:clearing}{clearing} arm. Overnight, the collateral that the US equity
\omterm{def:m1:clearing-and-settlement:clearing}{clearing} house required from it had risen to about \$3.7 billion. The firm
had about \$700 million on deposit; the remaining \$3 billion was due by
ten o'clock. Nothing had gone wrong with any trade. Its customers had bought
enormous quantities of a few violently moving stocks, those purchases would
not be paid for until two days later, and until then somebody had to stand
behind them. That somebody is the subject of this chapter: what happens
between the moment a trade is agreed and the moment it is final, who
guarantees it meanwhile, and what that guarantee costs.

\section{From trade to settlement}

\begin{definition}[Clearing and settlement]\label{def:m1:clearing-and-settlement:clearing}
\emph{Settlement}\index{settlement} is the final exchange of what was traded:
securities move to the buyer, cash to the seller, and the trade can no longer
be undone. \emph{Clearing}\index{clearing} is everything between execution
and settlement: confirming the terms, computing who owes what to whom,
netting, and managing the risk that a party fails before it has paid or
delivered.
\end{definition}

\begin{definition}[Settlement cycle and settlement fail]\label{def:m1:clearing-and-settlement:cycle}
The \emph{settlement cycle}\index{settlement cycle} is the standard delay
between trade date $T$ and \omterm{def:m1:clearing-and-settlement:clearing}{settlement} date, written $T{+}n$ for $n$ business
days. A \emph{settlement fail}\index{settlement fail} is an obligation not
met on its \omterm{def:m1:clearing-and-settlement:clearing}{settlement} date, usually because the seller did not have the
securities to deliver.
\end{definition}

\begin{dated}{2026-09}{Settlement cycles}\label{dat:m1:clearing-and-settlement:tplus1}
US securities have settled at $T{+}1$ since 28 May 2024. The European Union
has legislated the same move for 11 October 2027 (the amendment to its
central securities depositories regulation was published in October 2025);
the United Kingdom and Switzerland have aligned on the same date. Until then
most European securities settle at $T{+}2$.
\end{dated}

\begin{definition}[Central securities depository and delivery versus payment]\label{def:m1:clearing-and-settlement:csd}
A \emph{central securities depository}\index{central securities depository}
(CSD) keeps the definitive record of who owns each security of a market and
settles trades by book entry between its participants' accounts.
\emph{Delivery versus payment}\index{delivery versus payment} (DVP) is the
rule that the securities leg and the cash leg of a \omterm{def:m1:clearing-and-settlement:clearing}{settlement} become final
together or not at all, so that neither party can lose the full value of the
trade.
\end{definition}

\begin{omfigure}
\begin{tikzpicture}[font=\small,
  ev/.style={draw=omDef,rounded corners=2pt,minimum width=2.7cm,minimum height=0.85cm,align=center,fill=omDef!4}]
\draw[-{Stealth[length=2.2mm]},thick,omExo] (-0.6,0) -- (15.2,0);
\foreach \x/\lab in {0.6/{$T$, 10:31},4.4/{$T$, evening},8.2/{$T{+}1$, morning},12/{$T{+}1$, day}}{
  \draw[omExo,thick] (\x,-0.12) -- (\x,0.12);
  \node[font=\footnotesize,omExo,anchor=north] at (\x,-0.2) {\lab};
}
\node[ev] at (0.6,1.05) {Execution\\on a venue};
\node[ev] at (4.4,1.05) {Novation,\\netting};
\node[ev] at (8.2,1.05) {Margin call\\met};
\node[ev] at (12,1.05) {Settlement\\DVP at the CSD};
\draw[omThm,thick,decorate,decoration={brace,mirror,amplitude=5pt}] (0.6,-0.95) -- (12,-0.95)
  node[midway,below=6pt,font=\footnotesize,omThm] {the clearing house is exposed to each member's default};
\end{tikzpicture}

\omcaption{The life of a share trade under a one-day cycle. The brace is the
window the rest of the chapter is about: its length, times the volatility of
what was traded, sets the collateral.}\label{fig:m1:clearing-and-settlement:timeline}
\end{omfigure}

\section{The central counterparty}

\begin{definition}[Central counterparty and novation]\label{def:m1:clearing-and-settlement:ccp}
A \emph{central counterparty}\index{central counterparty} (CCP), or \omterm{def:m1:clearing-and-settlement:clearing}{clearing}
house, interposes itself in every trade it accepts: by
\emph{novation}\index{novation} the contract between buyer and seller is
replaced by two contracts, one between the buyer and the CCP and one between
the CCP and the seller. Each member then has a single counterparty, whatever
the number of firms it traded with.
\end{definition}

\begin{proposition}[What novation does to a network]\label{prop:m1:clearing-and-settlement:network}
Among $n$ members trading bilaterally there can be $n(n-1)/2$ counterparty
relationships; through a CCP there are $n$. After \omterm{def:m1:clearing-and-settlement:ccp}{novation} each member's
obligations in one security and one currency collapse to a single net
position against the CCP, and these net positions sum to zero.
\end{proposition}
\begin{proof}
A pair of members is a set of two among $n$: $\binom n2$ pairs. With a CCP
every contract has the CCP on one side, giving one relationship per member.
A member's contracts with the CCP in one security are fungible and add to
$\sum(\text{bought}) - \sum(\text{sold})$. Every trade adds $+q$ to one
member and $-q$ to another, so the sum over members is zero.
\end{proof}

\begin{omfigure}
\begin{tikzpicture}[font=\small,
  m/.style={circle,draw=omDef,fill=omDef!6,minimum size=0.75cm,inner sep=0pt},
  c/.style={circle,draw=omThm,fill=omThm!8,minimum size=1.05cm,inner sep=0pt}]
\begin{scope}
\foreach \i/\a in {A/90,B/30,C/330,D/270,E/210,F/150}{ \node[m] (\i) at (\a:1.9) {\i}; }
\foreach \i/\j in {A/B,A/C,A/D,A/E,A/F,B/C,B/D,B/E,B/F,C/D,C/E,C/F,D/E,D/F,E/F}{ \draw[omExo] (\i) -- (\j); }
\node[font=\footnotesize] at (0,-2.75) {bilateral: 15 relationships};
\end{scope}
\begin{scope}[xshift=7.6cm]
\node[c] (K) at (0,0) {CCP};
\foreach \i/\a in {A/90,B/30,C/330,D/270,E/210,F/150}{ \node[m] (\i) at (\a:1.9) {\i}; \draw[omDef,thick] (\i) -- (K); }
\node[font=\footnotesize] at (0,-2.75) {cleared: 6 relationships};
\end{scope}
\end{tikzpicture}

\omcaption{Six members before and after \omterm{def:m1:clearing-and-settlement:ccp}{novation}. The CCP removes the web,
and concentrates in one node all the risk that was spread over it.}
\end{omfigure}

\begin{example}[Netting one morning's trades]\label{ex:m1:clearing-and-settlement:netting}
A buys 5\,000 shares from B at \$20.00, B buys 3\,000 from C at \$20.10, and C
buys 4\,000 from A at \$19.95. Gross, 12\,000 shares and \$240\,100 would
move. Net, A receives 1\,000 shares and pays \$20\,200; B delivers 2\,000 and
receives \$39\,700; C receives 1\,000 and pays \$19\,500. Three deliveries
totalling 2\,000 shares and \$39\,700 settle instead: netting has removed
83\% of the value. The tutorial shows the ratio rising towards 98\% as
activity grows, which is why a market of thousands of members can settle at
all.
\end{example}

\section{Margin and the default waterfall}

A CCP that guarantees every trade must survive the default of a member. It
does so with other people's money, called in a fixed order.

\begin{definition}[Variation margin and initial margin]\label{def:m1:clearing-and-settlement:margin}
\emph{Variation margin}\index{variation margin} is the daily (or intraday)
payment of the change in value of a member's open positions: losers pay, the
CCP passes the cash to winners, and exposures restart from zero.
\emph{Initial margin}\index{initial margin} is collateral deposited against
the loss the CCP could suffer on the member's positions between its last
variation-margin payment and the moment the CCP has finished closing them
out, a delay called the margin period of risk.
\end{definition}

\begin{method}[A value-at-risk initial margin]\label{met:m1:clearing-and-settlement:im}
For a net position of value $N$ in an instrument of daily volatility
$\sigma$, to be covered over $d$ days with confidence $\Phi(z)$:
\[
\mathrm{IM} = z\,\sigma\sqrt{d}\;|N| .
\]
CCPs add charges for concentration, for illiquidity, for gaps between
volatility regimes and --- in the US equity \omterm{def:m1:clearing-and-settlement:clearing}{clearing} house --- for a member
whose requirement is large relative to its own capital. Real models are
portfolio-based; \cref{ch:m1:margin} builds one.
\end{method}

\begin{example}[Why the call came]\label{ex:m1:clearing-and-settlement:why}
Apply the method to a broker whose customers bought, net, \$1 billion of a
stock: with $\sigma = 4\%$, $d=2$, $z = 2.33$, $\mathrm{IM} = \$132$ million.
If the stock's volatility jumps to 25\% a day the same position requires
\$824 million, and doubling the position doubles it again. In January 2021
the volatility component of the retail broker's requirement was about \$1.3
billion, to which the \omterm{def:m1:clearing-and-settlement:clearing}{clearing} house's formula added an ``excess capital
premium'' of \$2.2 billion because the requirement dwarfed the firm's
capital. The \omterm{def:m1:clearing-and-settlement:clearing}{clearing} house waived that premium the same morning, for this
firm and others: \$9.7 billion in total across its members that week,
according to the congressional investigation. The broker restricted
purchases in the stocks concerned --- the only way it had to stop the
requirement from growing --- and set out the same week to raise \$3.5
billion of new capital. Shortening the cycle to $T{+}1$ divides such a
requirement by $\sqrt2$: that, not convenience, was the argument for it.
\end{example}

\begin{omfigure}
\begin{tikzpicture}
\begin{axis}[width=10.5cm,height=4.8cm,
  xlabel={daily volatility of the position (\%)},ylabel={margin (\% of position)},
  tick label style={font=\small},label style={font=\small},
  xmin=1,xmax=15,ymin=0,ymax=55,grid=major,grid style={gray!25},
  legend columns=2,legend style={font=\footnotesize,at={(0.5,-0.34)},anchor=north,draw=none,fill=none,/tikz/every even column/.append style={column sep=8pt}},
  table/col sep=comma]
\addplot[omThm,very thick] table[x=sigma_pct,y=t2]{figdata/markets-1/05-clearing-and-settlement/margin_vs_vol.csv};
\addplot[omDef,very thick] table[x=sigma_pct,y=t1]{figdata/markets-1/05-clearing-and-settlement/margin_vs_vol.csv};
\legend{two-day exposure,one-day exposure}
\end{axis}
\end{tikzpicture}

\omcaption{The value-at-risk margin of \cref{met:m1:clearing-and-settlement:im}
at 99\% confidence. Margin is linear in volatility: a stock that becomes
five times more volatile overnight costs five times more to clear the next
morning. Data: computed by the chapter's script.}
\end{omfigure}

\begin{definition}[Default fund and default waterfall]\label{def:m1:clearing-and-settlement:waterfall}
The \emph{default fund}\index{default fund} is a pool of collateral
contributed by all members to absorb losses that exceed a defaulter's own
resources. The \emph{default waterfall}\index{default waterfall} is the order
in which resources absorb the loss from a member's default:
\begin{enumerate}
\item the defaulter's \omterm{def:m1:clearing-and-settlement:margin}{initial margin};
\item the defaulter's contribution to the default fund;
\item a tranche of the CCP's own capital (its ``skin in the game'');
\item the surviving members' default-fund contributions;
\item further assessments the CCP may call from survivors, up to a cap.
\end{enumerate}
\end{definition}

The design is ``defaulter pays first, survivors pay next, together'': layers
1 and 2 make each member answer for its own risk; layer 3 gives the CCP's
owners a reason to set margins properly; layers 4 and 5 make every member a
guarantor of every other --- which is why members care who else is admitted
and how much they are allowed to do.

\begin{omfigure}
\begin{tikzpicture}
\begin{axis}[width=11cm,height=5.4cm,stack plots=y,area style,no markers,
  xlabel={loss from the default (\$ million)},ylabel={who pays (\$ million)},
  tick label style={font=\small},label style={font=\small},
  xmin=0,xmax=1100,ymin=0,ymax=1100,
  legend columns=3,legend style={font=\footnotesize,at={(0.5,-0.3)},anchor=north,draw=none,fill=none,/tikz/every even column/.append style={column sep=6pt}},
  table/col sep=comma]
\addplot[fill=omDef!70,draw=omDef] table[x=loss,y=margin]{figdata/markets-1/05-clearing-and-settlement/waterfall.csv} \closedcycle;
\addplot[fill=omDef!35,draw=omDef] table[x=loss,y=own_fund]{figdata/markets-1/05-clearing-and-settlement/waterfall.csv} \closedcycle;
\addplot[fill=omMeth!55,draw=omMeth] table[x=loss,y=ccp]{figdata/markets-1/05-clearing-and-settlement/waterfall.csv} \closedcycle;
\addplot[fill=omProp!55,draw=omProp] table[x=loss,y=survivors]{figdata/markets-1/05-clearing-and-settlement/waterfall.csv} \closedcycle;
\addplot[fill=omThm!45,draw=omThm] table[x=loss,y=assess]{figdata/markets-1/05-clearing-and-settlement/waterfall.csv} \closedcycle;
\addplot[fill=black!55,draw=black] table[x=loss,y=uncovered]{figdata/markets-1/05-clearing-and-settlement/waterfall.csv} \closedcycle;
\legend{defaulter's margin,defaulter's fund share,CCP capital,survivors' fund,assessments,uncovered}
\end{axis}
\end{tikzpicture}

\omcaption{Who absorbs a default loss as it grows, for illustrative layers of
120, 30, 20, 400 and 400 million. Up to 150 the defaulter pays for itself;
from 170 the other members do; beyond 970 nobody is committed to. Data:
computed by the chapter's script.}\label{fig:m1:clearing-and-settlement:waterfall}
\end{omfigure}

\section{When a member defaults}

\begin{method}[The default-management sequence]\label{met:m1:clearing-and-settlement:dm}
\begin{enumerate}
\item \textbf{Declare} the default and stop accepting the member's trades.
\item \textbf{Port} the defaulter's clients, with their positions and
  margin, to other members where the account structure allows it.
\item \textbf{Hedge} the defaulter's house portfolio to neutralise its market
  risk within hours.
\item \textbf{Auction} the hedged portfolio to the surviving members, who
  are obliged or strongly incentivised to bid.
\item \textbf{Allocate} any loss down the waterfall; return what remains of
  the defaulter's margin to its administrator.
\end{enumerate}
\end{method}

\begin{example}[September 2008]\label{ex:m1:clearing-and-settlement:lehman}
When Lehman Brothers defaulted on Monday 15 September 2008, the London
\omterm{def:m1:clearing-and-settlement:clearing}{clearing} house for interest-rate swaps held its portfolio: 66\,390 trades
with a notional value of \$9 trillion in five currencies, against about \$2
billion of \omterm{def:m1:clearing-and-settlement:margin}{initial margin}. Traders seconded from member banks hedged the
portfolio alongside the \omterm{def:m1:clearing-and-settlement:clearing}{clearing} house's risk team; between 24 September and
3 October the hedged currency portfolios were auctioned. The \omterm{def:m1:clearing-and-settlement:clearing}{clearing} house
reported that the default was managed well within the margin held and that
its \omterm{def:m1:clearing-and-settlement:waterfall}{default fund} was not used. The episode became the standard argument for
the \omterm{def:m1:clearing-and-settlement:clearing}{clearing} mandates that followed.
\end{example}

\begin{remark}[The risk has not disappeared]\label{rem:m1:clearing-and-settlement:concentration}
A CCP does not remove counterparty risk; it replaces many small, opaque
exposures with one large, collateralised and supervised one. Its margin
calls are \emph{procyclical}: they rise exactly when volatility is high and
cash is scarce, so that the mechanism protecting the system can drain the
liquidity of its members at the worst moment --- the theme of
\cref{ch:m1:margin,ch:m1:stress-case-studies}. And the \omterm{def:m1:clearing-and-settlement:clearing}{clearing} house itself
becomes a firm that cannot be allowed to fail.
\end{remark}

\section{Tutorial: netting a day of trades}\label{tut:m1:clearing-and-settlement:netting}

\begin{tutorial}
\textbf{Goal.} Net a day of trades multilaterally and measure what netting
removes. \textbf{End state:} the chart below, from 41\% of value removed for
ten trades to 98\% for ten thousand.
\begin{tutsteps}
\item \textbf{Net.} One pass over the trades accumulates, per member, shares
per symbol and cash.
\omcode{code/markets-1/05-clearing-and-settlement/python/clearing_demo.py}{18}{35}{Multilateral netting and the share of value it removes.}\label{lst:m1:clearing-and-settlement:net}
\item \textbf{Check the invariants.} Net cash sums to zero across members,
and so do net shares in each symbol: the tests assert both. A netting engine
that violates them has lost a trade.
\item \textbf{Measure.} Generate random trades among eight members in five
stocks and average the efficiency over twenty seeds.
\item \textbf{Allocate a loss.} The waterfall is six lines.
\omcode{code/markets-1/05-clearing-and-settlement/python/clearing_demo.py}{58}{65}{Walking a loss down the waterfall.}\label{lst:m1:clearing-and-settlement:wf}
\end{tutsteps}
\textbf{What to change next.} Make one member a retail broker that only buys
one stock: its own netting efficiency collapses, and with it the economy of
margin that netting gives everyone else. Then give each member an \omterm{def:m1:clearing-and-settlement:margin}{initial
margin} from \cref{met:m1:clearing-and-settlement:im} on its net position and
find the member whose requirement is largest relative to its size.
\end{tutorial}

\begin{omfigure}
\begin{tikzpicture}
\begin{axis}[width=10.5cm,height=4.6cm,xmode=log,
  xlabel={trades in the day (eight members, five stocks)},ylabel={value removed by netting (\%)},
  tick label style={font=\small},label style={font=\small},
  xmin=10,xmax=10000,ymin=30,ymax=100,grid=major,grid style={gray!25},table/col sep=comma]
\addplot[omDef,very thick,mark=*,mark size=1.4pt] table[x=trades,y=eff_pct]{figdata/markets-1/05-clearing-and-settlement/netting.csv};
\end{axis}
\end{tikzpicture}

\omcaption{Multilateral netting efficiency against activity (logarithmic
horizontal axis). The more two-way flow a market has, the smaller the
fraction that ever has to settle. Data: the tutorial's simulation, mean of
twenty seeds per point.}
\end{omfigure}

\section{Build: the netting engine}\label{bld:m1:clearing-and-settlement:netting}

\begin{build}
\bfield{Purpose} The miniature firm clears its own simulated market. At the
end of each session it must know, per member, what settles.
\bfield{Interface} \texttt{NettingEngine.add(trade)} during the session;
\texttt{close()} returning an immutable \texttt{SettlementInstructions}:
per member and symbol a signed share quantity, per member and currency a
signed cash amount in ledger units (\cref{ch:m1:what-a-trading-firm-does}),
and the list of trade identifiers netted into each.
\bfield{Rules} Prices and cash are integers in ledger units: no floating
point. A trade with buyer equal to seller, a non-positive quantity or a
duplicate identifier is rejected. After \texttt{close()} no trade is
accepted. The two zero-sum invariants are checked inside \texttt{close()}
and their violation is a fatal error, not a warning.
\bfield{Acceptance tests} \texttt{code/firm/clearing/tests/}: reproduces
\cref{ex:m1:clearing-and-settlement:netting} exactly; invariants on ten
thousand random trades; rejections; every trade identifier appears in exactly
two members' instructions.
\bfield{Stretch} Add \texttt{fail(member, symbol, quantity)}: a partial
delivery, the resulting fail carried to the next day, and the buy-in that
closes it.
\end{build}

\begin{omsources}
\item US House Committee on Financial Services, majority staff report,
\emph{Game Stopped: How the Meme Stock Market Event Exposed Troubling
Business Practices, Inadequate Risk Management, and the Need for Legislative
and Regulatory Reform}, June 2022.
\item US Securities and Exchange Commission, \emph{Staff Report on Equity and
Options Market Structure Conditions in Early 2021}, October 2021.
\item US Securities and Exchange Commission, press release 2024-62 on the
$T{+}1$ settlement cycle; Euronext, \emph{T+1 programme}; UK Financial
Conduct Authority, \emph{About T+1 settlement}.
\item LCH.Clearnet, \emph{\$9 trillion Lehman OTC interest rate swap default
successfully resolved}, press release, 8 October 2008.
\item Committee on Payments and Market Infrastructures and IOSCO,
\emph{Principles for financial market infrastructures}, April 2012.
\item J.~Gregory, \emph{Central Counterparties}, Wiley, 2014.
\end{omsources}

\section{Exercises}

\begin{exercise}[$\star$]\label{exo:m1:clearing-and-settlement:1}
How many bilateral relationships can exist among 40 members? How many with a
CCP? By what factor does the count fall?
\end{exercise}

\begin{exercise}[$\star$]\label{exo:m1:clearing-and-settlement:2}
A sells 2\,000 shares to B at \$50.00; B sells 2\,000 to C at \$50.20; C sells
1\,500 to A at \$50.10. Give each member's net shares and net cash, and the
netting efficiency by value.
\end{exercise}

\begin{exercise}[$\star$]\label{exo:m1:clearing-and-settlement:3}
Compute the \omterm{def:m1:clearing-and-settlement:margin}{initial margin} of \cref{met:m1:clearing-and-settlement:im} on a
net position of \$400 million with $\sigma = 2.5\%$, $z = 2.33$, for a
two-day and for a one-day period. What does the shorter cycle free up?
\end{exercise}

\begin{exercise}[$\star\star$]\label{exo:m1:clearing-and-settlement:4}
With the waterfall of \cref{fig:m1:clearing-and-settlement:waterfall},
allocate losses of (a) \$140 million, (b) \$300 million, (c) \$1\,000
million. There are 25 surviving members with equal fund contributions: what
does each lose in (b)?
\end{exercise}

\begin{exercise}[$\star\star$]\label{exo:m1:clearing-and-settlement:5}
A broker's customers hold a net unsettled long position of \$600 million.
Overnight the stock's daily volatility is re-estimated from 3\% to 18\%.
Compute the margin before and after ($z = 2.33$, two days). The broker's
capital is \$500 million: comment.
\end{exercise}

\begin{exercise}[$\star\star$]\label{exo:m1:clearing-and-settlement:6}
Without DVP, a seller delivers \$10 million of shares in the morning and is
to be paid in the afternoon. With DVP and a CCP, the trade is guaranteed at
its price. If the buyer fails at noon after the stock has fallen 4\%, what
does the seller lose in each case?
\end{exercise}

\begin{exercise}[$\star\star\star$]\label{exo:m1:clearing-and-settlement:7}
\emph{Coding.} Using \texttt{random\_trades(2000, 8, seed=1)}, compute the
netting efficiency. Then append 400 trades in which member \texttt{A} buys
500 shares of \texttt{S0} from a random other member, and recompute it for
the whole set. Explain the sign of the change.
\end{exercise}

\begin{exercise}[$\star\star\star$]\label{exo:m1:clearing-and-settlement:8}
\emph{Find the flaw.} ``Our \omterm{def:m1:clearing-and-settlement:clearing}{clearing} house has never used its \omterm{def:m1:clearing-and-settlement:waterfall}{default fund},
so its margins are clearly sufficient; we can reduce the fund and lower
members' costs.'' Give two reasons why the premise does not support the
conclusion, one statistical and one about behaviour.
\end{exercise}

\section{Problem: The Default of Member C}

\begin{problem}[{Weekend problem --- a clearing house on its worst Monday}]\label{pb:m1:clearing-and-settlement:1}
A \omterm{def:m1:clearing-and-settlement:clearing}{clearing} house for equity index futures has 21 members. Member C, a
proprietary firm, is long futures with a notional value of \$2.4 billion. The
\omterm{def:m1:clearing-and-settlement:clearing}{clearing} house holds from C an \omterm{def:m1:clearing-and-settlement:margin}{initial margin} of \$130 million and a
default-fund contribution of \$25 million. Its own capital tranche is \$15
million; the other 20 members contribute \$20 million each to the \omterm{def:m1:clearing-and-settlement:waterfall}{default
fund} and can each be assessed once more for the same amount.

\medskip\noindent\textbf{Part I --- Was the margin right?}
\begin{enumerate}
\item The index has a daily volatility of 1.3\%. What margin does
\cref{met:m1:clearing-and-settlement:im} give for $d = 2$ days at 99\%
($z = 2.33$)?
\item The \omterm{def:m1:clearing-and-settlement:clearing}{clearing} house held \$130 million. To what confidence level does
that correspond?
\item On Friday the index falls 3\%. What \omterm{def:m1:clearing-and-settlement:margin}{variation margin} is C called for?
\item C pays it. What is the \omterm{def:m1:clearing-and-settlement:clearing}{clearing} house's exposure to C on Friday
evening?
\end{enumerate}

\medskip\noindent\textbf{Part II --- The default.}
\begin{enumerate}[resume]
\item On Monday morning the index opens 6\% lower and C does not pay.
What does C owe?
\item The \omterm{def:m1:clearing-and-settlement:clearing}{clearing} house declares C in default at 9:00. Hedging the
portfolio takes until noon, during which the index falls another 2\%.
Compute the total loss on the position since Friday's close.
\item The auction of the hedged portfolio costs a further 0.4\% of notional
in concessions to the bidders. Give the total loss to be allocated.
\item Allocate it down the waterfall, layer by layer.
\item How much does each surviving member lose?
\end{enumerate}

\medskip\noindent\textbf{Part III --- Counterfactuals.}
\begin{enumerate}[resume]
\item If the \omterm{def:m1:clearing-and-settlement:clearing}{clearing} house had hedged within one hour, before the further
fall, what would the survivors have lost?
\item What \omterm{def:m1:clearing-and-settlement:margin}{initial margin} would have covered the whole loss of question 7?
Express it as a multiple of daily volatility times notional.
\item What is the largest loss the full waterfall can absorb?
\item To what total index fall since Friday's close does that correspond,
with the same auction cost?
\end{enumerate}

\medskip\noindent\textbf{Part IV --- Design.}
\begin{enumerate}[resume]
\item After the event the \omterm{def:m1:clearing-and-settlement:clearing}{clearing} house raises $z$ to 3 and the margin
period to three days. Give the new margin on a \$2.4 billion position.
\item C's position was 30\% of the open interest. Propose a
concentration add-on and justify its form.
\item Why does the \omterm{def:m1:clearing-and-settlement:clearing}{clearing} house's own capital sit before the survivors'
money and not after?
\item A member argues that assessments should be uncapped ``so that the
\omterm{def:m1:clearing-and-settlement:clearing}{clearing} house can never fail''. Give the argument against.
\item During the same Monday every other member received a large
intraday margin call. Why is this a risk to the system, and what is it
called?
\item State the \emph{named result}: the loss reaching the mutualised
\omterm{def:m1:clearing-and-settlement:waterfall}{default fund}, in millions of dollars.
\item In one sentence: who, in the end, guarantees a cleared trade?
\end{enumerate}
\end{problem}

\section{Interview questions}

\begin{interviewq}[$\star$ \iqroles{trader, developer, bank}]\label{iq:m1:clearing-and-settlement:1}
What is \omterm{def:m1:clearing-and-settlement:ccp}{novation}, and what risk does it remove for a trading firm?
\end{interviewq}

\begin{interviewq}[$\star$ \iqroles{trader, researcher, bank}]\label{iq:m1:clearing-and-settlement:2}
What is the difference between \omterm{def:m1:clearing-and-settlement:margin}{initial margin} and \omterm{def:m1:clearing-and-settlement:margin}{variation margin}? Which
one do you get back?
\end{interviewq}

\begin{interviewq}[$\star\star$ \iqroles{trader, researcher}]\label{iq:m1:clearing-and-settlement:3}
In January 2021 several retail brokers stopped customers from buying certain
stocks. Give the mechanical explanation.
\end{interviewq}

\begin{interviewq}[$\star\star$ \iqroles{bank, researcher}]\label{iq:m1:clearing-and-settlement:4}
Describe the \omterm{def:m1:clearing-and-settlement:waterfall}{default waterfall} of a \omterm{def:m1:clearing-and-settlement:clearing}{clearing} house. Why is the order what it
is?
\end{interviewq}

\begin{interviewq}[$\star\star$ \iqroles{developer}]\label{iq:m1:clearing-and-settlement:5}
You are writing the end-of-day netting job. What invariants do you assert
before publishing \omterm{def:m1:clearing-and-settlement:clearing}{settlement} instructions, and what do you do if one fails
at 18:55 with a 19:00 deadline?
\end{interviewq}

\begin{interviewq}[$\star\star\star$ \iqroles{researcher, bank}]\label{iq:m1:clearing-and-settlement:6}
Central \omterm{def:m1:clearing-and-settlement:clearing}{clearing} was made mandatory for standard swaps after 2008. Does it
reduce systemic risk? Argue both sides.
\end{interviewq}
